%%=============================================================================
%% Inleiding
%%=============================================================================

\chapter{\IfLanguageName{dutch}{Inleiding}{Introduction}}%
\label{ch:inleiding}

Deze bachelorproef richt zich op het automatiseren van CRA-compliance aan de hand van\\CI/CD-pipelines binnen de softwareontwikkeling bij Excentis. Door de invoering van de Europese Cyber Resilience Act (CRA) krijgen digitale producten striktere cybersecurityverplichtingen \parencite{EU_CRA}. Hierdoor moet software tegen december 2027 voldoen aan de vereisten van de CRA-wetgeving om een CE-label te verkrijgen \parencite{EuropeanCommission2024CRAnews}. Dit betekent dat de software die door Excentis wordt ontwikkeld conform aan deze wetgeving moet zijn. Daardoor vormt de nieuwe regelgeving een grote uitdaging voor zowel development teams als management teams. Volgens \textcite{Robyn2025} verwachten de klanten dat Excentis hier ook expliciet mee aan de slag gaat, waardoor het probleem een concrete en bedrijfsrelevante casus vormt. Zo wordt niet alleen CRA-compliance verzekerd, maar dit ondersteunt ook het klantenvertrouwen, waardoor de producten veilig en op tijd op de markt kunnen worden gebracht.
Het doel van deze bachelorproef is het ontwerpen en implementeren van een volledige CI/CD-pipeline die automatisch de CRA-compliance van de software controleert. Hierdoor zal het onderzoek ook een concrete oplossing bieden voor de specifieke probleemsituatie van Excentis, namelijk het waarborgen van veilige en CRA-compliant software, binnen de gestelde deadlines van de CRA. De gemaakte pipeline zal dienen als een Proof of Concept en laat Excentis toe om kwetsbaarheden te vinden, een Software Bill of Materials (SBOM) te genereren en hiertegen CVE-scans uit te voeren, waarbij er automatisch een rapportage plaatsvindt naar de betrokken development en management teams. 

\section{\IfLanguageName{dutch}{Probleemstelling}{Problem Statement}}%
\label{sec:probleemstelling}

De Cyber Resilience Act (CRA) verplicht producenten van digitale producten om hun software en hardware aantoonbaar veilig en betrouwbaar te maken vóór deze op de markt komen. Voor Excentis, dat software ontwikkelt, betekent dit dat elke update nog steeds aan de relevante CRA-vereisten voldoet. Dit vormt een grote uitdaging, omdat handmatige controles op beveiliging en compliance veel tijd en gespecialiseerde expertise vragen en daardoor moeilijk schaalbaar zijn. De bachelorproef richt zich daarom expliciet op Excentis als doelgroep en heeft tot doel een oplossing uit te werken die direct inzetbaar is binnen hun ontwikkelcontext.
\section{\IfLanguageName{dutch}{Onderzoeksvraag}{Research question}}%
\label{sec:onderzoeksvraag}

De centrale onderzoeksvraag is dan ook: 

Hoe kan een CI/CD-pipeline geautomatiseerd worden om CRA-compliant te zijn? 
Om deze onderzoeksvraag te beantwoorden, worden er verschillende deelvragen onderzocht binnen zowel het probleemdomein als het oplossingsdomein. Deelvragen in het probleemdomein zijn:
\begin{itemize}
  \item Wat is de Cyber Resilience Act (CRA), en welke doelstellingen en verplichtingen bevat het?
  \item Wat is CI/CD en waarom is dit belangrijk binnen DevOps?
  \item Wat is de impact van CRA voor Excentis, welke producten worden getroffen en wat zien zij als gewenste output?
\end{itemize}
Daarnaast worden ook de onderzoeksvragen uit het oplossingsdomein behandeld:
\begin{itemize}
  \item Hoe wordt een CI/CD-pipeline technisch opgebouwd binnen een DevOps-omgeving en welke voordelen biedt automatisatie voor beheer en beveiliging?
  \item Hoe kunnen beveiligingscontroles en compliancevereisten uit de Cyber Resilience Act geïntegreerd en geautomatiseerd worden binnen CI/CD-pipelines, welke CRA-vereiste securitycontroles daarbij geautomatiseerd moeten worden, en welke software hiervoor beschikbaar is, om bijvoorbeeld een SBOM te genereren?
  \item Welke technische en organisatorische uitdagingen kunnen optreden bij het automatiseren van CRA-compliance in CI/CD-pipelines?
\end{itemize}

\section{\IfLanguageName{dutch}{Onderzoeksdoelstelling}{Research objective}}%
\label{sec:onderzoeksdoelstelling}

De onderzoeksdoelstelling van deze bachelorproef is het ontwikkelen van een Proof of Concept van een CI/CD-pipeline die niet CRA-compliant software automatisch kan detecteren. Concreet moet de pipeline:
\begin{itemize}
  \item Automatisch de geselecteerde securitytesten uit een vooraf bepaalde shortlist uitvoeren
  \item Een SBOM genereren en deze vergelijken met een CVE-database
  \item De gevonden kwetsbaarheden en compliance-risico's rapporteren aan de betrokken development- en managementteams
\end{itemize}
De Proof of Concept wordt als geslaagd beschouwd wanneer de pipeline deze functionaliteiten correct en betrouwbaar uitvoert op representatieve voorbeeldprojecten en door Excentis als voldoende bruikbaar en uitbreidbaar wordt beoordeeld.

\section{\IfLanguageName{dutch}{Opzet van deze bachelorproef}{Structure of this bachelor thesis}}%
\label{sec:opzet-bachelorproef}

% Het is gebruikelijk aan het einde van de inleiding een overzicht te
% geven van de opbouw van de rest van de tekst. Deze sectie bevat al een aanzet
% die je kan aanvullen/aanpassen in functie van je eigen tekst.

De rest van deze bachelorproef is als volgt opgebouwd:

In Hoofdstuk~\ref{ch:stand-van-zaken} wordt een overzicht gegeven van de stand van zaken binnen het onderzoeksdomein, op basis van een literatuurstudie.

In Hoofdstuk~\ref{ch:methodologie} wordt de methodologie toegelicht en worden de gebruikte onderzoekstechnieken besproken om een antwoord te kunnen formuleren op de onderzoeksvragen.

% TODO: Vul hier aan voor je eigen hoofstukken, één of twee zinnen per hoofdstuk

In Hoofdstuk~\ref{ch:conclusie}, tenslotte, wordt de conclusie gegeven en een antwoord geformuleerd op de onderzoeksvragen. Daarbij wordt ook een aanzet gegeven voor toekomstig onderzoek binnen dit domein.